\originalchapter{网络流与割对偶}\label{ch:flow}
\sourceof{18.433 L7--8; 与 Menger、匹配的衔接为编者补充}
\originalsection{容量与守恒}
\begin{definition}
流网络是有限有向多重图, 指定不同顶点源 $s$、汇 $t$, 每条弧 $e$ 有容量 $c(e)\ge0$. 可行流为弧上的实数 $f(e)$, 满足 $0\le f(e)\le c(e)$, 且除 $s,t$ 外每点流出量等于流入量. 流值 $|f|$ 为源的净流出量, 也等于汇的净流入量.
\end{definition}
自环对流值无作用, 可删除. 平行弧与反平行弧各有独立编号. 对 $s\in S,t\notin S$ 的割, 容量为从 $S$ 到补集的弧容量之和, 反向弧不计入容量.
\begin{lemma}[弱对偶]\label{lem:flowweak}
任意可行流与任意 $s$--$t$ 割满足
\[
|f|=\sum_{e:S\to\overline S}f(e)-\sum_{e:\overline S\to S}f(e)\le c(S,\overline S).
\]
\end{lemma}
\begin{proof}
把 $S$ 内各顶点的净流出量求和, 内部弧一进一出抵消, 非源点由守恒贡献0, 所以留下源的流值. 前向流不超过容量, 反向流非负, 得不等式.
\end{proof}
\originalsection{残量网络与增广}
对每条原弧 $e=u\to v$, 建立前向残量弧 $u\to v$ 容量 $c(e)-f(e)$ 和后向残量弧 $v\to u$ 容量 $f(e)$, 只保留正残量弧. 它们保留原边编号与方向身份, 因而反平行原弧不会与另一弧的回退混淆. 前向表示还能增加流量, 后向表示可以撤回已安排流量.

若残量网络有 $s$--$t$ 路, 令瓶颈 $\theta$ 为路上最小残量. 沿前向原弧加 $\theta$, 沿后向减 $\theta$. 简单路不会同时使用同一原弧两种身份. 更新后仍可行, 中间点净变化为0, 源净流出增加 $\theta$. 无后向弧的“只向前加流”容易陷入局部选择, 不能作为一般最大流算法.
\begin{theorem}[最大流最小割]\label{thm:maxflow}
有限网络的最大流值等于最小割容量. 若所有容量为非负整数, 则有整数最大流, 从零流反复沿增广路更新可在有限次后得到它.
\end{theorem}
\begin{proof}
整数容量时, 每次残量和瓶颈都是整数, 一次至少使流值增加1. 流值由源的出弧总容量界住, 所以有限次后无增广路. 令 $S$ 为残量中从源可达点集, 汇不在其中. 每条原前向跨割弧已饱和, 否则还能到达其终点; 每条原反向跨割弧流量为0, 否则其后向残量会从 $S$ 到补集. 所以弱对偶式取等, 流值等于该割容量, 两者都最优.

一般实容量时, 不保证任意选增广路的过程终止. 但可行流构成有限维闭有界集合, 非空, 线性流值在其上达到最大. 最大流若残量仍有路, 正瓶颈会改进流值, 矛盾. 因而它无增广路, 按相同可达集构造割就取等. 这里的存在性仅用有限维紧致性, 不依赖 Menger 或匹配定理.
\end{proof}
\originalsection{Edmonds--Karp 的时间界}
每次在残量网络用 BFS 选择边数最少的增广路, 就得到 Edmonds--Karp 算法. 这不仅选出一条容易找到的路, 还阻止某条弧不断在同一距离层上重复成为瓶颈.
\begin{lemma}\label{lem:residualdistance}
采用最短增广路时, 各顶点从源的残量距离不会减小, 不可达点距离记为无限.
\end{lemma}
\begin{proof}
更新前各旧残量弧 $u\to v$ 满足 $d(v)\le d(u)+1$. 更新后新出现的弧只可能是增广路边的反向. 对增广路上的 $u\to v$, 有 $d(v)=d(u)+1$, 因而新反向弧也满足 $d(u)\le d(v)+1$. 新图中若有比旧距离短的路, 从源沿它逐边应用这些不等式, 该路长仍至少为旧终点距离, 矛盾. 旧时不可达点无法由新反向弧首次到达, 因增广路各点旧时已可达, 所以无限距离也不会下降.
\end{proof}
\begin{theorem}
Edmonds--Karp 至多增广 $O(nm)$ 次, 用邻接表与 BFS 的时间为 $O(nm(n+m))$, 在连通可达网络的通常 $m\ge n-1$ 情形可写为 $O(nm^2)$.
\end{theorem}
\begin{proof}
一次增广至少有一个残量弧饱和, 称为本次临界弧. 若 $u\to v$ 是临界弧, 当时 $d(v)=d(u)+1$. 要使同一身份的残量弧再次出现并再次饱和, 中间必须沿其反向 $v\to u$ 增广以恢复残量. 在那次最短路上, 新距离满足 $d'(u)=d'(v)+1\ge d(v)+1=d(u)+2$, 用了距离不减引理. 所以同一残量弧两次临界之间, 尾点距离至少增加2. 有限可达距离至多 $n-1$, 每种残量身份最多临界 $O(n)$ 次, 总身份数至多 $2m$, 得增广次数. 每次 BFS 扫描 $O(n+m)$ 项, 得时间界. 删除与源汇问题无关的孤立顶点后, 常用界为 $O(nm^2)$.
\end{proof}
证明没有要求容量为整数, 所以算法对精确实数运算也有限终止. 实际计算机处理浮点残量需明确容差; 这里是数学算法时间界, 不把近似浮点比较当作精确证书.
\originalsection{一个流与割的证书}
\begin{example}
网络有弧 $sa:3,sb:2,ab:1,at:2,bt:3$. 求最大流及最小割.
\end{example}
\begin{solution}
先沿 $s,a,t$ 送2, 再沿 $s,b,t$ 送2, 最后沿 $s,a,b,t$ 送1. 得 $f(sa)=3,f(sb)=2,f(ab)=1,f(at)=2,f(bt)=3$, 各中间点守恒, 总流5. 割 $S=\{s\}$ 容量 $3+2=5$, 流值与割容量相等, 所以无需再搜索也能验证最优.
\end{solution}
\begin{center}
\begin{tikzpicture}[x=2.1cm,y=1.15cm]
\node[vertex](s)at(0,0){$s$};\node[vertex](a)at(1,1){$a$};\node[vertex](b)at(1,-1){$b$};\node[vertex](t)at(2,0){$t$};
\draw[arc](s)--node[above left]{$3/3$}(a);\draw[arc](s)--node[below left]{$2/2$}(b);
\draw[arc](a)--node[right]{$1/1$}(b);\draw[arc](a)--node[above right]{$2/2$}(t);\draw[arc](b)--node[below right]{$3/3$}(t);
\end{tikzpicture}
\end{center}
标注为“流量/容量”. 容量上界与守恒是两套约束, 校验证书时二者都必须检查.
\originalsection{匹配、路与对偶的统一}
二部图 $A\sqcup B$ 的最大匹配可化为流: 加源到每个 $A$ 点容量1, 每条二部边作为 $A\to B$ 容量1, 每个 $B$ 点到汇容量1. 整数流使每点最多参与一条正流二部边, 所以流值就是匹配大小; 反过来每个匹配给等值流. 顶点容量拆分则给出第\ref{ch:connectivity}章 Menger 的顶点版本. 同一个守恒与残量机制因此联系了边不相交路、顶点不相交路与匹配.

最大流还可写为线性规划, 最小割为一个特殊的对偶证书. 本册使用残量可达集实际构造这个证书, 不先假设线性规划强对偶. 一般费用流、单纯形、匹配多面体与花算法的复杂度不在本册范围.
\originalsection{习题与解答}
\begin{exercise}\label{ex:flow1}
网络弧为 $sa,sb,ac,ad,bc,ct,dt$, 容量全为1. 若先沿 $s,a,c,t$ 增广, 给出第二次增广路以达到最大流2.
\end{exercise}
\begin{exercise}\label{ex:flow2}
设 $f$ 是一个整数最大流, 说明每个整数 $0\le k\le |f|$ 也能作为某个整数流的值.
\end{exercise}
习题\ref{ex:flow1}的解答.
\begin{solution}
残量路 $s,b,c,a,d,t$ 使用后向弧 $c\to a$, 撤回第一次的 $a\to c$, 改成两条路 $s,b,c,t$ 与 $s,a,d,t$. 流值2, 源总容量2给最优上界. 如果不允许回退, 第一次选择会错误地看似只能送1.
\end{solution}
习题\ref{ex:flow2}的解答.
\begin{solution}
把整数流的圈流先消去, 剩余正流反复沿源汇路抽取一单位, 得 $|f|$ 条单位路流. 取其中任意 $k$ 条相加, 仍满足原容量与守恒, 得整数值 $k$. 这只是可行流的分解, 不要求单位路边不相交, 因为容量可以大于1.
\end{solution}
